% ---------------------------------------------------------------------
%  5 UŽDUOTIS — Eksperimentinis vertinimas
%  REIKALINGA PREAMBULĖJE: xltabular, tabularx, booktabs, siunitx,
%                          graphicx, float, multirow
%  Lentelės ir paveikslai GENERUOJAMI:
%    python -m src.eksperimentai.i_latex --aibe test --priesaga _test
%    python -m src.eksperimentai.klaidos --aibe test
%    python -m src.eksperimentai.slenkstis --taikyti test
%    python -m src.eksperimentai.nematytos
%    python -m src.eksperimentai.patikimumas --su-duomenimis
%    python -m src.eksperimentai.paveikslai
% ---------------------------------------------------------------------

\subsection{Vertinimo protokolas ir operacinis taškas}
\label{sec:vertinimo_protokolas}

Testavimo aibė iki šio skyriaus nebuvo atidaryta nė karto. Visi
sprendimai --- hiperparametrai, ankstyvas stabdymas, požymių sąrašas ir
sprendimo slenkstis --- priimti mokymo ir validacijos aibėse, o testavimo
aibė panaudota vienu prėjimu, iš kurio gauti visi šiame skyriuje
pateikiami pjūviai.

Vertinami keturi modeliai suderintomis konfigūracijomis, po tris
paleidimus su skirtingais pradiniais dydžiais; pateikiamas vidurkis ir
standartinis nuokrypis. Pagrindinė formuluotė --- aštuonios kategorijos;
dvejetainė ir 34 klasių nagrinėjamos atskirai (\ref{sec:formuluotes}
poskyris).

Palyginimas atliekamas ties \textbf{suderintu klaidingų teigiamų
biudžetu}, o ne ties didžiausios tikimybės taisykle. Slenkstis
$\tau$ kiekvienam modeliui parinktas validacijos aibėje kaip mažiausias,
tenkinantis 1~\% biudžetą, ir testavimo aibėje netaikomas iš
naujo: perrinkimas toje pačioje aibėje, kurioje matuojama, būtų
nutekėjimas. Didžiausios tikimybės taškas pateikiamas šalia, nes ties juo
modelių rikiuotė yra kitokia --- tai savarankiškas rezultatas, aptariamas
\ref{sec:slenkscio_kompromisas} poskyryje.

Skirtumas tarp modelių laikomas tikru tik tada, kai jis viršija
paleidimų sklaidą. Formalus statistinis testas neatliekamas: trijų
paleidimų $t$~testas turi pervertintą pirmos rūšies klaidos tikimybę
\cite{dietterich1998tests}.


\subsection{Pagrindiniai rezultatai}
\label{sec:pagrindiniai_rezultatai}

\begin{table}[htbp]
\centering
\caption{Aptikimo kokybė testavimo aibėje.}
\label{tab:rezultatai_test}
\lentele{rezultatai_test.tex}
\end{table}

Ties didžiausios tikimybės tašku geriausią makro-F1 pasiekia
Random~Forest (\num{0.7276}), po jo XGBoost (\num{0.7210}) ir MLP
(\num{0.6346}). Skirtumas tarp dviejų pirmųjų yra \num{0.0066}, o
paleidimų sklaida --- \num{0.0001}, todėl jis nėra atsitiktinis.

Ties klaidingų teigiamų biudžetu rikiuotė apsiverčia: XGBoost pasiekia
makro-F1 \num{0.6626} prie 0,95~\% klaidingų teigiamų, o
Random~Forest --- \num{0.6455} prie 1,02~\%. Priežastis yra
tikimybių skiriamoji geba: Random~Forest tikimybė yra medžių balsų dalis,
todėl ties vienetu ji kinta šuoliais, o XGBoost grąžina tolydžias
reikšmes ir aukštą slenkstį pasiekia tiksliau. Tas pats reiškinys buvo
išmatuotas ir validacijos aibėje, ir su nesuderintomis
hiperparametrų reikšmėmis, todėl tai modelių savybė, o ne vieno matavimo
ypatybė.

Autokoderis makro-F1 skalėje su prižiūrimais modeliais nelyginamas: jis
sprendžia dvejetainį uždavinį, o jo \num{0.2187} gaunamas prie
0,93~\% klaidingų teigiamų --- taško, kurį jis pasiekia pagal
konstrukciją, nes slenkstis yra gerybinio srauto atkūrimo paklaidos
procentilis.

\begin{table}[htbp]
\centering
\caption{Veikimo rodikliai testavimo aibėje: mokymo laikas, gryna
  inferencijos delsa procesoriumi ir modelio dydis diske.}
\label{tab:veikimas_test}
\lentele{veikimas_test.tex}
\end{table}

Inferencijos delsa visiems keturiems modeliams yra
3--30 mikrosekundžių vienam įrašui, t.~y. tris eiles mažesnė už
20--50~ms kraštinio šliuzo biudžetą, nustatytą
\ref{sec:atakos} skyriuje. Ribojantis resursų rodiklis yra ne delsa, o
\textbf{modelio dydis}: suderintas Random~Forest užima 638~MB
diske, o įkeliamas į atmintį --- kelis kartus daugiau, todėl kraštinio
šliuzo klasės įrenginyje, turinčiame 1--8~GB, jis netelpa. XGBoost su
45~MB telpa be išlygų.


\subsection{Klaidų analizė}
\label{sec:klaidu_analize}

Bendras tikslumas visas klaidas skaičiuoja vienodai, nors eksploatacijai
jos kainuoja skirtingai. \ref{tab:klaidu_tipai} lentelėje klaidos
suskirstytos į tris rūšis: klaidingi teigiami (gerybinis srautas
paskelbtas ataka), praleistos atakos ir painiava \emph{tarp atakų}, kai
pavojaus signalas įvyksta, o suklystama tik kategorija.

\begin{table}[htbp]
\centering
\caption{Klaidų pasiskirstymas pagal eksploatacinę kainą.}
\label{tab:klaidu_tipai}
\lentele{klaidu_tipai.tex}
\end{table}

Beveik keturios klaidos iš penkių yra tarp atakų. Iš XGBoost
\num{61651} klaidos \num{39155} sudaro viena pora: \num{33379} DDoS
eilučių priskirtos DoS ir \num{5776} DoS eilučių --- DDoS, t.~y.
63,5~\% visų klaidų. Abi kategorijos kraštiniame šliuze
sukelia tą patį veiksmą, o jas skirianti informacija --- srauto šaltinių
skaičius ir kryptis --- šiame duomenų rinkinio leidime neišreikšta.

Iš to seka, kad tikslumas \num{0.8303} nereiškia, jog modelis praleidžia
kas šeštą ataką: du trečdaliai atotrūkio nuo vieneto yra riba, kurios
turimi požymiai neleidžia nubrėžti, o ne neaptiktos atakos.

\begin{figure}[H]
  \centering
  \includegraphics[width=0.82\textwidth]{sumaisymas.pdf}
  \caption{XGBoost sumaišymo matrica testavimo aibėje, normalizuota per
    eilutę. Ryškiausi ne įstrižainės langeliai --- DDoS ir DoS riba bei
    gerybinio srauto ir žvalgybos pora.}
  \label{pav:sumaisymas}
\end{figure}

Antra stambi klaidų grupė turi vardą, numatytą dar
\ref{sec:atakos} skyriuje. Tikras gerybinis srautas, kurį modelis
paskelbia ataka, beveik visas virsta \emph{žvalgyba}: Random~Forest ---
21,2~\% viso gerybinio srauto, XGBoost --- 14,5~\%.
Painiava abipusė: \num{7770} žvalgybos eilučių priskiriamos gerybiniam
srautui, todėl gerybinio srauto tikslumas tesiekia \num{0.543}.

Tai tiesioginė \ref{tab:atakos} lentelėje užfiksuoto apribojimo pasekmė.
Žvalgybos eilutėje ten nurodyta, kad srauto kryptis, trukmė ir unikalių
taikinių skaičius šiame leidime nefiksuojami --- o būtent jie skiria
prievadų skenavimą nuo įprasto jungčių srauto. Teorinėje dalyje įvardytas
požymių trūkumas eksperimente pasirodė kaip konkretus, išmatuojamas klaidų
šablonas.

\begin{table}[htbp]
\centering
\caption{Per-kategorijų rodikliai testavimo aibėje.}
\label{tab:perklase}
\lentele{perklase.tex}
\end{table}

Makro-F1 lygį lemia dvi mažiausios kategorijos. \emph{Web}
($n = \num{3556}$) ir \emph{BruteForce} ($n = \num{1878}$) F1 yra
0,38--0,47 medžių ansambliuose ir 0,21--0,22 MLP,
o \emph{Mirai} --- \num{0.998}. Retose kategorijose didžiausia ir
paleidimų sklaida, todėl pagrindinio rodiklio neapibrėžtumas kyla būtent
iš jų.

MLP klysta kitaip nei medžių ansambliai: jo klaidingi teigiami eina ne į
žvalgybą, o į \emph{Web} ir \emph{BruteForce}, kurių tikslumas tesiekia
0,125--0,137 prie atkūrimo 0,60--0,62. Klasių
svoriai jį verčia retas kategorijas skelbti per dažnai, o medžių
ansambliuose tas pats svertas veikia priešinga kryptimi.


\subsection{Klaidingi teigiami ir slenksčio kompromisas}
\label{sec:slenkscio_kompromisas}

\begin{table}[htbp]
\centering
\caption{Operaciniai taškai testavimo aibėje. $\tau$ parinktas
  validacijos aibėje.}
\label{tab:slenkstis_test}
\lentele{slenkstis_test.tex}
\end{table}

Ties didžiausios tikimybės tašku prižiūrimų modelių klaidingų teigiamų
dalis yra 21--31~\% --- 21--31 kartą daugiau už
\ref{sec:atakos} skyriuje nustatytą biudžetą. Įvedus slenkstį, XGBoost
klaidingi teigiami krenta iki 0,95~\%, o aptikimas išlieka
88,5~\%: klaidingų teigiamų sumažėja daugiau nei
dvidešimt kartų, o makro-F1 --- 8~\%. Modelis nepermokomas;
keičiasi tik sprendimo taisyklė.

\begin{figure}[H]
  \centering
  \includegraphics[width=\textwidth]{kreives.pdf}
  \caption{Klaidingų teigiamų ir aptikimo kompromisas validacijos aibėje.
    Kairėje --- visi keturi modeliai; dešinėje priartinti prižiūrimieji su
    pažymėtais operaciniais taškais.}
  \label{pav:kreives}
\end{figure}

Autokoderio kreivė yra kitos rūšies. Tarp 95-ojo ir 99-ojo procentilio
jo aptikimas krenta nuo 80,9~\% iki 11,8~\%, o
prižiūrimi modeliai tame pačiame rėžyje lieka virš 83~\%.
Autokoderis \emph{rikiuoja} gerai --- jo PR-AUC didesnis už atsitiktinio
rikiavimo lygį --- bet reikalaujamo klaidingų teigiamų biudžeto viduje
naudingo aptikimo lygio nepasiekia. Tai stipresnis teiginys nei
konstatavimas, kad jis silpnesnis: jis pasako, kuriame taške ir kodėl.

\begin{figure}[H]
  \centering
  \includegraphics[width=\textwidth]{kategorijos.pdf}
  \caption{Aptikimas pagal kategorijas ties klaidingų teigiamų biudžetu
    testavimo aibėje. Gerybinio srauto eilutės nėra: jai tas pats
    stulpelis reikštų klaidingus teigiamus, ne aptikimą.}
  \label{pav:kategorijos}
\end{figure}

Bendras aptikimo rodiklis yra svertinis vidurkis, kurį lemia gausiausios
kategorijos. Ties tuo pačiu operaciniu tašku DDoS, DoS ir Mirai
aptinkamos 100~\%, o \emph{Recon} --- 43,9~\%,
\emph{Web} --- 37,5~\%. Rėžis per kategorijas yra
37--100~\%, ir prasčiausiai aptinkamos yra kaip tik tos,
kurios saugumo požiūriu svarbios kaip ankstyvas įspėjimas. Klaidingų
teigiamų biudžeto laikymasis perkamas žvalgybos aptikimu.


\subsection{Nematytų atakų klasių testas}
\label{sec:nematytos}

Neprižiūrimo metodo įtraukimas į palyginimą rėmėsi prielaida, kad
nematytoms atakoms reikia modelio, kuris atakų pavyzdžių nemato. Ši
prielaida patikrinta tiesiogiai: prižiūrimas modelis permokytas be
kiekvienos iš trijų klasių, parinktų dar \ref{sec:parinkimas} skyriuje,
ir įvertintas ta pačia testavimo aibe.

Klausiama vieno dalyko --- ar pašalintos klasės eilutės pažymimos kaip
\emph{bet kuri} ataka. Makro-F1 čia netinka: pašalinus
\texttt{DICTIONARYBRUTEFORCE} ištuštėja visa \emph{BruteForce}
kategorija, ir modelis vidurkinamas per septynias klases vietoj aštuonių.
Autokoderis nepermokomas --- jis mokomas tik iš gerybinio srauto, todėl
visos 33 atakų klasės jam ir taip nematytos.

\begin{table}[htbp]
\centering
\caption{Nematytų atakų klasių aptikimas testavimo aibėje.}
\label{tab:nematytos}
\lentele{nematytos.tex}
\end{table}

Prižiūrimas modelis, klasės niekada nematęs, aptinka ją geriau už
neprižiūrimą dviem atvejais iš trijų, o trečiuoju jie lygūs
(24,4~\% prieš 25,4~\%). Prielaida, kad zero-day
atakoms būtinas neprižiūrimas metodas, šiuose duomenyse nepasitvirtina.

Svarbesnis už rikiuotę yra klasės nematymo \emph{kaina}, matoma
sugretinus paskutinius du stulpelius:

\begin{itemize}\itemsep2pt
  \item \texttt{DDOS-SLOWLORIS} --- 0,0 procentinio punkto;
  \item \texttt{RECON-PORTSCAN} --- 1,5 procentinio punkto;
  \item \texttt{DICTIONARYBRUTEFORCE} --- 21,3 procentinio punkto.
\end{itemize}

Skirtumą lemia ne pati klasė, o tai, ar išlieka jos kategorija. Pirmųjų
dviejų kategorijose lieka atitinkamai vienuolika ir keturios giminingos
klasės, ir apibendrinimas joms beveik nemokamas. Trečioji yra vienintelė
savo kategorijos klasė, todėl ją pašalinus prižiūrimas modelis netenka
visos kategorijos, ir aptikimas krenta perpus.

\texttt{DDOS-SLOWLORIS} aptinkama 99,94~\% --- tiek pat, kiek
turint ją mokymo aibėje. Tai neprieštarauja \ref{sec:atakos} skyriuje
nurodytam apribojimui, o iš jo seka: kadangi srauto trukmės požymio
šiame leidime nėra, žemo intensyvumo ataka požymių erdvėje nesiskiria nuo
kitų DDoS klasių. Ji aptinkama patikimai ne todėl, kad modelis ją
atpažįsta, o todėl, kad neatskiria nuo to, kas jau išmokta ---
eksploatacijai signalas įvyksta, bet kategorija nurodoma neteisingai.

Teiginio ribos aiškios: tikrintos trys klasės, kurių kategorijos, išskyrus
vieną, mokymo aibėje liko. Rezultatas galioja naujai tos pačios šeimos
klasei, o ne atakai, nepriklausančiai nė vienai iš aštuonių kategorijų;
tokio atvejo šie duomenys neturi.


\subsection{Dvejetainė ir 34 klasių formuluotės}
\label{sec:formuluotes}

Palyginimui su literatūra geriausias prižiūrimas modelis įvertintas ir
kitomis dviem formuluotėmis, tais pačiais hiperparametrais. Palyginimas
atliekamas ties didžiausios tikimybės tašku, nes tokius skaičius skelbia
cituojami darbai.

\begin{table}[htbp]
\centering
\small
\begin{tabularx}{\textwidth}{@{}>{\raggedright\arraybackslash}X
    >{\centering\arraybackslash}p{2.0cm}>{\centering\arraybackslash}p{2.0cm}
    >{\centering\arraybackslash}p{2.0cm}>{\centering\arraybackslash}p{2.0cm}@{}}
\toprule
\textbf{Formuluotė} & \multicolumn{2}{c}{\textbf{Šis darbas}}
  & \multicolumn{2}{c}{\textbf{\cite{almahaqeri2026gradient}}} \\
\cmidrule(lr){2-3}\cmidrule(l){4-5}
 & tikslumas & makro-F1 & tikslumas & makro-F1 \\
\midrule
Dvejetainė      & 0,9435 & 0,7692 & 0,9961 & 0,9952 \\
8 kategorijos   & 0,8303 & 0,7210 & 0,9959 & 0,8903 \\
34 klasės       & 0,7708 & 0,6464 & 0,9948 & 0,8876 \\
\midrule
\textbf{Rėžis}  & \textbf{0,173} & 0,123 & \textbf{0,001} & 0,108 \\
\bottomrule
\end{tabularx}
\caption{Užduoties detalumo įtaka. Paskutinė eilutė --- skirtumas tarp
  didžiausios ir mažiausios reikšmės stulpelyje.}
\label{tab:formuluotes}
\end{table}

Skiriasi ne tik lygis, bet ir dėsningumo forma. Cituojamame darbe
užduoties detalumas tikslumui beveik neturi įtakos --- rėžis
\num{0.001} --- o makro-F1 krenta \num{0.108}. Šiame darbe makro-F1
kritimas panašus (\num{0.123}), o tikslumo rėžis yra
\num{0.173}, t.~y. daugiau nei šimtą kartų didesnis.

Toks skirtumas yra tai, ko tikėtumeisi, jei mokymo aibėje liktų
pasikartojančių eilučių. Įsiminta eilutė atsakoma teisingai bet kokiu
detalumu --- jai nesvarbu, ar klasių dvi, aštuonios, ar 34. Duomenyse,
iš kurių pašalinta 53,3~\% tikslių dublikatų, modelis privalo
klases tikrai atskirti, todėl kiekvienas detalumo lygis kainuoja.
Sugretinimą riboja ir kiti skirtumai --- 39 požymiai vietoj 46 bei imties
dydis --- todėl teigiama tiek, kad skiriasi dėsningumo forma, o ne tik
lygis.

Didžiausias atotrūkis yra paprasčiausioje užduotyje. Dvejetainėje
formuluotėje viskas priklauso nuo gerybinio srauto ir atakos ribos ---
tos pačios, kurioje glūdi \ref{sec:klaidu_analize} poskyryje aprašyta
žvalgybos painiava, --- o daugiaklasėje makro-F1 tą ribą atskiedžia
lengvai atpažįstamos kategorijos. Dvejetainis tikslumas \num{0.9435}
atrodo aukštas tik todėl, kad gerybinis srautas sudaro
4,1~\% eilučių.

Detalumas kainuoja pagal visas ašis. Perėjus nuo aštuonių kategorijų prie
34 klasių, modelis padidėja nuo 45~MB iki
113~MB, mokymas pailgėja nuo 112~s iki
407~s, delsa auga nuo 29,6 mikrosekundės iki
90,1 mikrosekundės, o makro-F1 krenta. Priežastis struktūrinė:
gradientinio stiprinimo modelis daugiaklasėje užduotyje augina
$n_{\text{medžių}} \times n_{\text{klasių}}$ medžių. Aštuonių kategorijų
detalumas, pasirinktas \ref{sec:di_metodai} skyriuje, taip pasirodo esąs
pagrįstas ir matavimu.


\subsection{Rezultatų patikimumas}
\label{sec:patikimumas}

Prieš rašant skaičius atliktos penkios patikros. Trys jų galioja be
išlygų, dvi davė radinius, kurie yra rezultato dalis.

\textbf{Testavimo aibė neįtakojo nė vieno sprendimo.} Slenkstis $\tau$
testavimo aibėje sutampa su validacijos aibėje parinktu iki paskutinio
skaitmens visiems keturiems modeliams, o kiekviena testavimo aibės eilutė
turi atitinkamą validacijos aibės eilutę --- vadinasi, nė vienas modelis
nebuvo vertintas testavimo aibėje anksčiau nei validacijos.

\textbf{Didžiausias tikslumas neviršija teorinės ribos.} Darbinėje imtyje
0,43~\% eilučių turi prieštaringas etiketes, todėl geriausias
įmanomas klasifikatorius klysta bent tiek; riba yra
99,78~\%, o didžiausias išmatuotas tikslumas ---
\num{0.9436}. Aukštesnis už ribą rezultatas būtų nutekėjimo požymis, ne
pasiekimas.

\textbf{Lentelių skaičiai atkuriami nepriklausomu keliu.} Kokybės
metrikos perskaičiuotos iš išsaugotų sumaišymo matricų ir sugretintos su
eksperimentų failu: didžiausias skirtumas dvylikoje paleidimų ---
$4{,}5 \cdot 10^{-6}$.

\begin{table}[htbp]
\centering
\caption{Operacinio taško perkeliamumas iš validacijos į testavimo aibę.}
\label{tab:patikimumas}
\lentele{patikimumas.tex}
\end{table}

\textbf{Pirmasis radinys --- operacinio taško perkeliamumas.} Klaidingų
teigiamų dalis testavimo aibėje nuo validacijos aibės vertės skiriasi
0,95--1,13 karto, tad taškas persikelia. Tačiau Random~Forest
peržengia biudžetą: 1,02~\% vietoj 1~\%.
Priežastis yra parinkimo taisyklėje --- imamas \emph{mažiausias} $\tau$,
tenkinantis biudžetą, todėl jis pagal konstrukciją atsiduria prie pat
ribos ir atsargos neturi. Eksploatacijai iš to seka konkretus
reikalavimas: slenkstį rinkti su atsarga, pavyzdžiui, ties
80~\% biudžeto.

\textbf{Antrasis radinys --- dedublikavimo ir požymių atrankos tvarka.}
Mokymo ir testavimo aibės nesikerta 39 stulpelių erdvėje, kurioje buvo
šalinami dublikatai. Tačiau 36 požymių erdvėje, kurią mato modelis,
sutampa \num{30} eilučių, t.~y. 0,0082~\% testavimo aibės.
Priežastis --- pašalinti požymiai. Eilutės, kurios skiriasi paskutiniu
slankiojo kablelio bitu vien pašalintame stulpelyje, po požymių atrankos
tampa tapačios. Poveikis rezultatams neviršija \num{0.0082} procentinio
punkto, t.~y. lieka žemiau pateikiamo tikslumo, tačiau iš to seka
bendresnė taisyklė: dedublikavimas ir požymių atranka turi vykti toje
pačioje erdvėje.


\subsection{Apibendrinimas}
\label{sec:vertinimo_apibendrinimas}

Testavimo aibėje išmatuoti visi keturi modeliai trimis paleidimais,
atlikta klaidų analizė, nematytų atakų klasių testas ir palyginimas su
literatūra, o rezultatai patikrinti penkiomis patikromis.

Geriausias sprendimas ties reikalaujamu klaidingų teigiamų biudžetu yra
XGBoost: makro-F1 \num{0.6626}, aptinkama 88,5~\% atakų prie
0,95~\% klaidingų teigiamų, modelio dydis
45~MB, delsa 29,6 mikrosekundės. Dviejose užduoties
vietose rezultatą riboja ne modelis, o duomenys: gerybinio srauto ir
žvalgybos riba bei DoS ir DDoS skirtis, kurioms reikalingų krypties ir
trukmės požymių šiame leidime nėra.

\ref{sec:di_metodai} skyriuje nurodytas realistinis makro-F1 taikinys
0,85--0,90 nepasiektas. Trys išmatuoti veiksniai paaiškina
skirtumą: pašalinti 53,3~\% tikslių dublikatų, 39 požymių
leidimas vietoj 46 ir palyginimas ties klaidingų teigiamų biudžetu, o ne
ties didžiausios tikimybės tašku. Nė vienas iš jų nėra prielaida.
